% Options for packages loaded elsewhere
\PassOptionsToPackage{unicode}{hyperref}
\PassOptionsToPackage{hyphens}{url}
\PassOptionsToPackage{dvipsnames,svgnames,x11names}{xcolor}
\documentclass[
  10pt,
  fontset=fandol]{ctexart}
\usepackage{xcolor}
\usepackage[margin=16mm]{geometry}
\usepackage{amsmath,amssymb}
\setcounter{secnumdepth}{-\maxdimen} % remove section numbering
\usepackage{iftex}
\ifPDFTeX
  \usepackage[T1]{fontenc}
  \usepackage[utf8]{inputenc}
  \usepackage{textcomp} % provide euro and other symbols
\else % if luatex or xetex
  \usepackage{unicode-math} % this also loads fontspec
  \defaultfontfeatures{Scale=MatchLowercase}
  \defaultfontfeatures[\rmfamily]{Ligatures=TeX,Scale=1}
\fi
\usepackage{lmodern}
\ifPDFTeX\else
  % xetex/luatex font selection
\fi
% Use upquote if available, for straight quotes in verbatim environments
\IfFileExists{upquote.sty}{\usepackage{upquote}}{}
\IfFileExists{microtype.sty}{% use microtype if available
  \usepackage[]{microtype}
  \UseMicrotypeSet[protrusion]{basicmath} % disable protrusion for tt fonts
}{}
\makeatletter
\@ifundefined{KOMAClassName}{% if non-KOMA class
  \IfFileExists{parskip.sty}{%
    \usepackage{parskip}
  }{% else
    \setlength{\parindent}{0pt}
    \setlength{\parskip}{6pt plus 2pt minus 1pt}}
}{% if KOMA class
  \KOMAoptions{parskip=half}}
\makeatother
\setlength{\emergencystretch}{3em} % prevent overfull lines
\providecommand{\tightlist}{%
  \setlength{\itemsep}{0pt}\setlength{\parskip}{0pt}}
\usepackage{amsmath,amssymb,mathtools,needspace}
\xeCJKDeclareCharClass{CJK}{"2032,"0400->"04FF,"2160->"216B,"2460->"2473,"25A0,"25C7,"25CB,"25CF}
\IfFontExistsTF{Arial Unicode MS}{\xeCJKsetup{AutoFallBack=true}\setCJKfallbackfamilyfont{\CJKrmdefault}{Arial Unicode MS}}{}
\setcounter{secnumdepth}{0}
\setlength{\emergencystretch}{2em}
\allowdisplaybreaks
\usepackage{bookmark}
\IfFileExists{xurl.sty}{\usepackage{xurl}}{} % add URL line breaks if available
\urlstyle{same}
\hypersetup{
  pdftitle={直接法的缩减技术},
  colorlinks=true,
  linkcolor={Maroon},
  filecolor={Maroon},
  citecolor={Blue},
  urlcolor={Blue},
  pdfcreator={LaTeX via pandoc}}

\title{直接法的缩减技术}
\author{}
\date{}

\begin{document}
\maketitle

\subsection{13.2 直接法的缩减技术}\label{ux76f4ux63a5ux6cd5ux7684ux7f29ux51cfux6280ux672f}

\textbf{所谓直接法是这样一类算法，它通过有限步计算可以直接得出问题的精确解}（如果不考虑舍入误差的话）。

\subsubsection{13.2.1 数列求和的累加算法}\label{ux6570ux5217ux6c42ux548cux7684ux7d2fux52a0ux7b97ux6cd5}

下述数列求和问题是人们所熟知的：

\[
S=a_0+a_1+\cdots+a_n.
\tag{13.2.1}
\]

这个计算模型有两个简单的特例。当 \(n=0\) 即为一项和式 \(S=a_0\) 时，所给计算模型就是它的解，这时不需要做任何计算。这表明，对于数列求和问题，它的解是计算模型退化的情形。又当 \(n=1\) 即计算两项和式 \(S=a_0+a_1\) 时，计算过程是平凡的，这时不存在算法设计问题。

现在基于这两种简单情形考察所给和式（13.2.1）的累加求和算法。设 \(b_k\) 表示前 \(k+1\) 项的部分和 \(a_0+a_1+\cdots+a_k\)，则有

\href{../../source-images/pdf-309.jpeg}{原页图像}

\[
\begin{cases}
b_0=a_0,\\
b_k=b_{k-1}+a_k,\quad k=1,2,\cdots,n,
\end{cases}
\tag{13.2.2}
\]

而计算结果 \(b_n\) 即为所求的和值

\[
S=b_n.
\tag{13.2.3}
\]

上述\textbf{数列求和的累加算法，其设计思想是将多项求和（式（13.2.1））化归为两项求和（式（13.2.2））的重复。}而依式（13.2.2）重复加工若干次，最终即可将所给和式（13.2.1）加工成一项和式（13.2.3）的退化情形，从而得出和值 \(S\)。

再剖析计算模型自身的演变过程。按式（13.2.2）每加工一次，所给和式（13.2.1）便减少一项，而所生成的计算模型依然是数列求和。反复施行这种加工手续，计算模型不断变形为

\[
\underset{\text{（计算模型）}}{n+1\text{ 项和式}}
\Rightarrow n\text{ 项和式}
\Rightarrow n-1\text{ 项和式}
\Rightarrow\cdots\Rightarrow
\underset{\text{（所求结果）}}{1\text{ 项和式}},
\]

这里，符号``\(\Rightarrow\)''表示重复施行两项求和的加工手续。

这样，如果定义和式的项数为数列求和问题的\textbf{规模}，则所求和值可以视作规模为 1 的退化情形。因此，只要令和式的规模（项数）逐次减 1，最终当规模为 1 时即可直接得出所求的和值。这样设计出的算法就是累加求和算法。

上述累加求和算法可以视作规模缩减技术的一个范例。

\subsubsection{13.2.2 缩减技术的设计机理}\label{ux7f29ux51cfux6280ux672fux7684ux8bbeux8ba1ux673aux7406}

许多数值计算问题可以引进某个实数------所谓问题的\textbf{规模}来刻画其``大小''，而问题的解则是其规模为足够小，譬如规模为 1 或 0 的退化情形。求解这类问题，一种行之有效的办法是通过某种简单的运算手续逐步缩减问题的规模，直到加工得出所求的解。算法设计的这种技术称作\textbf{规模缩减技术}，简称\textbf{缩减技术}。

\textbf{缩减技术适用的一类问题是，求解这类问题的困难在于它的规模（适当定义）比较大。针对这类问题运用缩减技术，就是设法逐步缩减计算问题的规模，直到规模变得足够小时直接生成或方便地求出问题的解。}

缩减技术的设计机理可用\textbf{``大事化小，小事化了''}这句俗话来概括。

\textbf{所谓``大事化小''意即逐步压缩问题的规模。}在运用缩减技术时，``大事''是如何``化小''的呢？这个处理过程具有如下两项基本特征。

（1）\textbf{结构递归}：``大事化小''是逐步完成的，其每一步将所考察的计算模型加工成同样类型的计算模型，因而这类算法具有明晰的递归结构。

（2）\textbf{规模递减}：每一步加工前后的计算模型虽然从属于同一类型，但其规模已被压缩了。压缩系数愈小，算法的效率愈高。

再考察``小事化了''的处理过程。\textbf{所谓``小事化了''，是指当问题的规模变得足够小时即可直接或方便地得出问题的解。}

\href{../../source-images/pdf-310.jpeg}{原页图像}

``小事''是如何``化了''的呢？

对于某些计算模型，如前面讨论过的数列求和问题，它们的规模为正整数，而其解则是规模为 0 或 1 的退化情形。这时只要设法使规模逐次减 1，加工若干步后即可直接得出所求的解。这里``小事化了''是直截了当的。

这样设计出的一类算法统称\textbf{直接法}。前述数列求和的累加算法以及下面将要讲到的多项式求值的秦九韶算法都是直接法。

\subsubsection{13.2.3 多项式求值的秦九韶算法}\label{ux591aux9879ux5f0fux6c42ux503cux7684ux79e6ux4e5dux97f6ux7b97ux6cd5}

微积分方法的核心是逼近法。多项式是微积分学中最为基本的一种逼近工具，因而多项式求值算法在微积分计算中具有重要意义。

设要对给定的 \(x\) 计算下列多项式的值：

\[
P=a_0x^n+a_1x^{n-1}+\cdots+a_{n-1}x+a_n
=\sum_{k=0}^{n}a_kx^{n-k}.
\tag{13.2.4}
\]

由于计算每一项 \(a_kx^{n-k}\) 需做 \(n-k\) 次乘法，如果先逐项计算 \(a_kx^{n-k}\)，然后再累加求和得出多项式的值 \(P\)，这种\textbf{逐项生成算法}所要耗费的乘法次数为

\[
Q=\sum_{k=0}^{n}(n-k)\approx\frac{n^2}{2}.
\]

当 \(n\) 充分大时，这个计算量是相当大的。

现在设法改进这一算法。类似于数列求和计算，首先考察两个特例：当 \(n=0\) 时，所给计算模型即为所求的解

\[
P=a_0,
\]

这时不需要做任何计算；又当 \(n=1\) 时，计算模型

\[
P=a_0x+a_1
\]

为简单的一次式，这时虽然需要进行计算，但不存在算法设计问题。

注意，当 \(x=1\) 时多项式（13.2.4）便退化为和式（13.2.1），可以类比数列求和算法的设计过程讨论多项式求值算法的设计问题。

设将多项式的次数规定为多项式求值问题的规模，如果从式（13.2.4）的前两项中提出公因子 \(x^{n-1}\)，则有

\[
P=(a_0x+a_1)x^{n-1}+\sum_{k=2}^{n}a_kx^{n-k}.
\]

这样，如果算出一次式

\[
v_1=a_0x+a_1
\]

的值，则所给 \(n\) 次式的计算模型（13.2.4）便化归为 \(n-1\) 次式

\[
P=v_1x^{n-1}+\sum_{k=2}^{n}a_kx^{n-k}
\]

\href{../../source-images/pdf-311.jpeg}{原页图像}

的计算，从而使问题的规模减少了 1 次。不断地重复这种加工手续，使计算问题的规模逐次减 1，则经过 \(n\) 步即可将所给多项式的次数降为 0，从而获得所求的解：

\textbf{算法 13.1}（秦九韶算法）　令 \(v_0=a_0\)，对 \(k=1,2,\cdots,n\) 计算

\[
v_k=xv_{k-1}+a_k,
\tag{13.2.5}
\]

则结果 \(P=v_n\) 即为所给多项式（13.2.4）的值。

容易看出，按递推算式（13.2.5）计算多项式（13.2.4）的值，总共只要做 \(n\) 次乘法，其计算量远比前述逐项生成算法的计算量小。这是一种优秀算法。

这一优秀算法称作\textbf{秦九韶算法}。它是公元 13 世纪我国南宋大数学家秦九韶（1208---1261）最先提出来的。需要注意的是，国外文献常称这一算法为 Horner 算法，其实 Horner 的工作比秦九韶晚了五六百年。

\textbf{秦九韶算法说明，\(n\) 次式（13.2.4）的求值问题可化归为一次式（13.2.5）求值计算的重复。}设以符号``\(\Rightarrow\)''表示一次式的求值手续，则秦九韶算法的模型加工流程如下：

\[
\underset{\text{（计算模型）}}{n\text{ 次式求值}}
\Rightarrow n-1\text{ 次式求值}
\Rightarrow n-2\text{ 次式求值}
\Rightarrow\cdots\Rightarrow
\underset{\text{（计算结果）}}{0\text{ 次式求值}}.
\]

\end{document}
